\subsection*{Spectral benchmark on $S^3$ and Hopf-compatible diagnostics}
  \label{subsec:hopf_spectral}

  The ratio $\lambda_2/\lambda_1=8/3$ is an exact spectral property of the scalar Laplace--Beltrami operator on the
  unit three-sphere.
  The Hopf fibration $S^1\hookrightarrow S^3\to S^2$ does not determine this ratio.
  Its role in the present appendix is operational: it provides a canonical fiber/base partition used to define the
  kernel-weighted moment observable $R(\alpha)$ in Eq.~\eqref{eq:R_alpha_def}, which is then shown to approach the same
  benchmark value in a Hopf-compatible projectable regime.

  \paragraph{1. Exact spectral ratio of $\Delta_{S^3}$.}
    For the unit-radius three-sphere and the positive Laplace--Beltrami operator, the scalar spectrum is
    \begin{equation}
      \lambda_k = k(k+2), \qquad k=0,1,2,\ldots,
    \end{equation}
    with multiplicity $(k+1)^2$.
    Hence $\lambda_1=3$ and $\lambda_2=8$, so that
    \begin{equation}
      \frac{\lambda_2}{\lambda_1} = \frac{8}{3}.
    \end{equation}
    This ratio is independent of any fibration structure and serves as the exact continuous benchmark established in
    Appendix~\ref{sec:analytical_ratio}.

  \paragraph{2. Geometric splitting induced by the Hopf fibration.}
    The Hopf fibration induces a canonical decomposition of the tangent bundle into a vertical (fiber) line and a
    horizontal (base) plane.
    This motivates the separation of squared increments into $d_{\mathrm{fiber}}^2$ and $d_{\mathrm{base}}^2$ used in the
    anisotropic kernel $K_\alpha$ of Eq.~\eqref{eq:anisotropic_kernel}.
    This geometric partition is used to probe how spectral responses separate along fiber-sensitive versus base-sensitive
    variations, but it does not alter the spectrum of $\Delta_{S^3}$ itself.

  \paragraph{3. Eigenspace structure and interpretation.}
    The first non-trivial eigenspace ($k=1$) has dimension $(1+1)^2=4$.
    Once a Hopf structure is fixed, it is convenient to interpret this space in terms of a singlet plus a triplet under
    the $\mathrm{SO}(3)$ symmetry acting on the base degrees of freedom.
    This $1+3$ decomposition is an interpretative tool for organizing modes relative to the fiber/base partition.
    It should not be confused with the tangent-space splitting $1+2$ or with any statement about the eigenvalue ratio.

  \paragraph{4. Connection to the weighted observable.}
    The convergence $R(\alpha)\to 8/3$ reported in Fig.~\ref{fig:convergence-8-3-s3} should be read as follows.
    The kernel $K_\alpha$ of Eq.~\eqref{eq:anisotropic_kernel} biases the edge ensemble so that fiber-sensitive increments
    contribute more strongly as $\alpha$ increases.
    In that Hopf-compatible regime, the normalized response defined in Eq.~\eqref{eq:normalized_ratio} approaches the exact
    spectral benchmark $8/3$.
    The kernel does not ``create'' the value $8/3$.
    It reveals the benchmark by isolating the fiber/base partition in the measured moments and by suppressing
    non-projectable contributions.

  \paragraph{5. Summary.}
    The ratio $\lambda_2/\lambda_1=8/3$ is an exact spectral invariant of the unit three-sphere.
    The Hopf fibration provides a canonical geometric partition used to define fiber/base diagnostics on graphs sampling
    $S^3$.
    The observable $R(\alpha)$ is a kernel-weighted moment ratio that approaches the same benchmark in a Hopf-compatible
    projectable regime, demonstrating operational robustness under refinement without asserting an operator-level identity
    between $R(\alpha)$ and $\lambda_2/\lambda_1$.
